\chapter{Feature Engineering and Selection}
\label{ch:feature-engineering}

\epigraph{``The single most important thing you can do to improve your model is to create better features.''}{}

\learninggoal{
	\item Understand why feature engineering is often the highest-leverage activity in applied ML.
	\item Apply common feature transformations: scaling, encoding, binning, interaction terms.
	\item Handle text, categorical, temporal, and missing data effectively.
	\item Select features using filter, wrapper, and embedded methods.
	\item Interpret feature importance using SHAP and permutation methods.
}

\section{Why Feature Engineering Matters}

Features are the input representation for a model. The quality of this representation determines the ceiling on model performance. A linear model with excellent features can outperform a neural network with raw inputs.

\begin{principle}[Data-Centric AI]
	Improving data quality and feature engineering often yields larger gains than improving model architecture or hyperparameter tuning~\cite{ng2024machine}.
\end{principle}

\section{Feature Transformation}

\subsection{Scaling}

Many models are sensitive to feature scales:

\begin{itemize}
	\item \textbf{Standardization (Z-score):} $x' = \frac{x - \mu}{\sigma}$. Centres at 0, scales to unit variance.
	\item \textbf{Min--Max scaling:} $x' = \frac{x - \min(x)}{\max(x) - \min(x)}$. Maps to $[0, 1]$.
	\item \textbf{Robust scaling:} $x' = \frac{x - \text{median}(x)}{\text{IQR}(x)}$. Resistant to outliers.
\end{itemize}

Decision trees and tree ensembles are scale-invariant. Linear models, SVMs, and neural networks are not.

\subsection{Log Transformation}

For right-skewed distributions (income, population, counts):

\[
	x' = \log(x + 1)
\]

This compresses the dynamic range and can make relationships more linear.

\subsection{Binning (Discretization)}

Converting a continuous feature into discrete bins can capture non-linear effects:

\begin{lstlisting}[language=Python, caption=Equal-width binning]
import pandas as pd
pd.cut(age, bins=[0, 18, 35, 60, 100],
       labels=['child', 'young', 'adult', 'senior'])
\end{lstlisting}

\begin{itemize}
	\item \textbf{Equal-width:} bins of equal range.
	\item \textbf{Equal-frequency (quantile):} bins with equal number of samples.
	\item \textbf{Adaptive:} bin boundaries chosen by domain knowledge or tree splits.
\end{itemize}

\subsection{Interaction Features}

When the effect of one feature depends on another, include their product:

\[
	x_3 = x_1 \times x_2
\]

Including all pairwise interactions expands $d$ features to $d + d(d-1)/2$. Tree-based models discover interactions automatically; linear models require explicit construction.

\section{Encoding Categorical Features}

\subsection{One-Hot Encoding}

Create a binary column for each category:

\[
	\text{color} = \{\text{red}, \text{green}, \text{blue}\} \to
	\begin{bmatrix}
		1 & 0 & 0 \\
		0 & 1 & 0 \\
		0 & 0 & 1
	\end{bmatrix}
\]

Dummy encoding drops one column $(K-1$ columns for $K$ categories) to avoid perfect multicollinearity.

\subsection{Target Encoding}

Replace each category with the mean target value for that category:

\[
	x_{\text{cat}}' = \frac{1}{n_{\text{cat}}} \sum_{i \in \text{cat}} y_i
\]

This is powerful but prone to overfitting, especially for rare categories. Smooth with a prior:

\[
	x_{\text{cat}}' = \frac{n_{\text{cat}} \bar{y}_{\text{cat}} + \lambda \bar{y}_{\text{global}}}{n_{\text{cat}} + \lambda}
\]

\subsection{Frequency Encoding}

Replace category with its count in the training set. Useful when frequency correlates with the target.

\section{Text Features}

\subsection{TF--IDF}

Term Frequency--Inverse Document Frequency transforms text into numerical vectors:

\[
	\text{TF--IDF}(t, d) = \text{TF}(t, d) \times \log\frac{N}{\text{DF}(t)}
\]

\begin{itemize}
	\item TF$(t, d)$: frequency of term $t$ in document $d$.
	\item DF$(t)$: number of documents containing term $t$.
	\item $N$: total documents.
\end{itemize}

TF--IDF down-weights common words (``the'', ``and'') and up-weights rare, informative terms.

\subsection{N-grams and Hashing}

\begin{itemize}
	\item \textbf{Unigrams:} individual words.
	\item \textbf{Bigrams:} pairs of consecutive words.
	\item \textbf{Feature hashing:} map arbitrary features to a fixed-size vector using a hash function (no dictionary needed).
\end{itemize}

\section{Temporal and Date Features}

Raw timestamps are rarely useful. Engineer:

\begin{itemize}
	\item Hour of day, day of week, month, quarter.
	\item Is weekend? Is holiday?
	\item Time since last event, rolling statistics (7-day mean, 30-day sum).
	\item Cyclical encoding: $\sin(2\pi \cdot \text{hour} / 24)$, $\cos(2\pi \cdot \text{hour} / 24)$.
\end{itemize}

\section{Handling Missing Data}

Strategies, in order of preference:

\begin{enumerate}
	\item \textbf{Remove:} drop rows or columns with excessive missingness ($> 50\%$).
	\item \textbf{Mean/median imputation:} fill with the feature mean. Simple but naive.
	\item \textbf{Model-based imputation:} predict missing values from observed ones (MICE, KNN imputer).
	\item \textbf{Indicator:} add a binary column ``is missing''.
	\item \textbf{Distributional:} sample from the observed distribution of the feature.
\end{enumerate}

\begin{principle}[Missing Not at Random]
	If the fact that a value is missing is informative about the target (e.g., survey non-response correlates with income), then proper treatment of missingness is critical. An indicator feature captures this signal.
\end{principle}

\section{Feature Selection}

Feature selection reduces dimensionality, improves generalization, and speeds up training.

\subsection{Filter Methods}

Rank features independently of the model:

\begin{itemize}
	\item \textbf{Correlation:} $|\rho(X_j, y)|$ for regression; ANOVA $F$-statistic for classification.
	\item \textbf{Mutual information:} $I(X_j; Y) = \sum_{x_j} \sum_y p(x_j, y) \log \frac{p(x_j, y)}{p(x_j) p(y)}$
	\item \textbf{Variance threshold:} remove features with near-zero variance.
\end{itemize}

\subsection{Wrapper Methods}

Evaluate feature subsets by training a model on each:

\begin{itemize}
	\item \textbf{Forward selection:} start empty, add features one by one.
	\item \textbf{Backward elimination:} start with all features, remove the least important.
	\item \textbf{Recursive feature elimination (RFE):} train, rank, remove bottom features, repeat.
\end{itemize}

Wrapper methods are computationally expensive but account for feature interactions.

\subsection{Embedded Methods}

Feature selection integrated into model training:

\begin{itemize}
	\item \textbf{Lasso (L1 regularization):} drives irrelevant coefficients to zero~\cite{tibshirani1996regression}.
	\item \textbf{Tree importance:} feature importance from tree splits.
	\item \textbf{Random forest importance:} mean decrease in impurity.
\end{itemize}

\section{Feature Importance and Interpretability}

\subsection{Permutation Importance}

Randomly shuffle a feature's values and measure the drop in performance:

\[
	\text{Importance}(X_j) = \text{Score}(\text{original}) - \text{Score}(\text{shuffled } X_j)
\]

This measures how much the model relies on $X_j$, regardless of the model type.

\subsection{SHAP Values}

SHAP (SHapley Additive exPlanations)~\cite{lundberg2017unified} decomposes a prediction into additive feature contributions based on cooperative game theory~\cite{shapley1953value}:

\[
	f(\mathbf{x}) = \phi_0 + \sum_{j=1}^d \phi_j
\]

\begin{itemize}
	\item $\phi_0$: baseline prediction (mean over training data).
	\item $\phi_j$: contribution of feature $j$ for this specific prediction.
	\item SHAP values satisfy desirable properties: local accuracy, missingness, consistency.
\end{itemize}

SHAP provides both global feature importance (mean $|\phi_j|$ across all data) and local explanations (why a specific prediction was made).

\section{The Feature Engineering Pipeline}

\begin{lstlisting}[language=Python, caption=Feature engineering pipeline]
from sklearn.pipeline import Pipeline
from sklearn.compose import ColumnTransformer
from sklearn.preprocessing import StandardScaler, OneHotEncoder
from sklearn.ensemble import RandomForestRegressor

numeric_features = ['age', 'income', 'hours']
categorical_features = ['education', 'occupation']

preprocessor = ColumnTransformer([
    ('num', StandardScaler(), numeric_features),
    ('cat', OneHotEncoder(), categorical_features)
])

pipeline = Pipeline([
    ('preprocess', preprocessor),
    ('model', RandomForestRegressor())
])

pipeline.fit(X_train, y_train)
\end{lstlisting}

\section{Exercises}

\begin{exercise}
	Take a continuous feature with a right-skewed distribution (e.g., income). Compare a model with the raw feature vs.\ $\log(\text{income})$. How do the residuals change? Which version produces more interpretable coefficients?
\end{exercise}

\begin{exercise}
	Apply one-hot encoding and target encoding to the same categorical feature with 50 categories. Train a linear model and a tree-based model. Compare the results. Why does target encoding generally help linear models more than tree models?
\end{exercise}

\begin{exercise}
	Given 100 features and 500 samples, use Lasso regression to select features. Vary the regularization parameter $\lambda$ and plot the number of non-zero coefficients vs.\ $\lambda$. How does cross-validation choose $\lambda$?
\end{exercise}

\begin{exercise}
	Using SHAP, explain why a specific test point received its prediction. List the top 3 features and their contributions. Are these globally important features or locally specific?
\end{exercise}
